% !PRINTED 456
\weberpage{456}

%<<CONT
Daraus ergiebt sich:

\[
\begin{aligned}
\tfrac{1}{6} F''(y_1, y_1) &= x_1 x_3 \alpha^2 + x_2 x_1 \beta^2 + x_3 x_2 \gamma^2 \\
&\quad + x_3^2 \beta\gamma + x_1^2 \alpha\gamma + x_2^2 \alpha\beta, \\
\tfrac{1}{3} F'''(y_2, y_3) &= 2x_1 x_3 \beta\gamma + 2x_2 x_1 \alpha\gamma + 2x_3 x_2 \alpha\beta \\
&\quad + x_3^2 (\beta\gamma + \alpha^2) + x_1^2 (\alpha\gamma + \beta^2) + x_2^2 (\alpha\beta + \gamma^2).
\end{aligned}
\]
% [sic] mixed term printed with triple prime F''' here against F'' just below for the same pairing

Da nun die Auflösungen des Systems (3) von derselben Form sind ($y_1 = \alpha x_1 + \beta x_2 + \gamma x_3, \dots$), so erhält man hieraus:

\[
\begin{aligned}
\tfrac{1}{6} F''(y_1, y_1) - y_1 y_3 &= \\
x_1 x_3 (\alpha^2 - \alpha\beta - \gamma^2) + x_2 x_1 (\beta^2 - \beta\gamma - \alpha^2) &+ x_2 x_3 (\gamma^2 - \beta^2 - \alpha\gamma), \\
\tfrac{1}{3} F''(y_2, y_3) - y_3^2 &= \\
x_1^2 (\beta^2 + \alpha\gamma - \gamma^2) + x_2^2 (\gamma^2 + \alpha\beta - \alpha^2) &+ x_3^2 (\alpha^2 + \beta\gamma - \beta^2).
\end{aligned}
\]

Die Coëfficienten auf der rechten Seite dieser Gleichungen $(\alpha^2 - \alpha\beta - \gamma^2) \dots$ ergeben sich aber aus den Werthen von $\alpha, \beta, \gamma$ [\S.\,115, (11)] als verschwindend, und wir erhalten also, wenn wir noch eine cyklische Vertauschung der $x$, die eine cyklische Vertauschung der $y$ zur Folge hat, anwenden:

\[
\begin{aligned}
&\tfrac{1}{6} F'''(y_1, y_1) = y_1 y_3, \quad \tfrac{1}{6} F'''(y_2, y_2) = y_2 y_1, \quad \tfrac{1}{6} F'''(y_3, y_3) = y_3 y_2, \\
&\tfrac{1}{3} F''(y_2, y_3) = y_3^2, \quad \tfrac{1}{3} F''(y_3, y_1) = y_1^2, \quad \tfrac{1}{3} F''(y_1, y_2) = y_2^2.
\end{aligned}
\]
% [sic] prime-count on F wavers between the two lines above, reproduced as printed

Demnach ergiebt sich nach dem \gesperrt{Euler}'schen Satze [Bd.\ I, \S.\,17, (6)]:

\[
F\,(y_1, y_2, y_3) = y_1^3 y_3 + y_2^3 y_1 + y_3^3 y_2,
\]

und damit ist nachgewiesen, dass die Function $f\,(x_1, x_2, x_3)$ in der That eine Invariante unserer Gruppe $G_{168}$ ist. Es ist, abgesehen von einem willkürlich beizufügenden constanten Factor, die einzige Invariante $4^{\mathrm{ten}}$ Grades.

Die Gleichung

\weq{5}{f\,(x_1, x_2, x_3) = 0}

bedeutet, wenn $x_1, x_2, x_3$ Coordinaten in der Ebene sind, eine \gesperrt{Curve vierter Ordnung}, die wir die \gesperrt{Grundcurve} der Gruppe oder die Curve $f$ nennen wollen, und es giebt 168 Abbildungen der Ebene auf sich selbst, bei denen jedem Punkte dieser Curve ein Punkt der Curve entspricht. Wenn ein Punkt auf der Grundcurve liegt, so liegen auch alle seine Bildpunkte darauf. Die Anzahl der verschiedenen Bildpunkte eines Punktes kann nur dann auf eine kleinere Zahl als 168, und zwar immer auf einen Theiler von 168, heruntersinken, wenn der Punkt ein Pol ist.
